\section{Capítulo~\ref{sec:glutcomputer}. O que os neurônios \nt{GLUT} calculam}

As afirmações lógicas do capítulo são teoremas sobre circuitos de pesos não negativos, deduzidos no próprio capítulo; a experiência confirma o modelo de neurônio e o fato de que circuitos montados segundo esses teoremas (detector de coincidência, linha de atraso, grupo autossustentado, cadeia) de fato existem no cérebro.

{\footnotesize
\setlength{\tabcolsep}{3pt}\renewcommand{\arraystretch}{1.15}
\begin{longtable}{>{\raggedright}p{0.5cm}>{\raggedright}p{4.0cm}>{\raggedright}p{5.6cm}>{\raggedright}p{2.3cm}>{\raggedright\arraybackslash}p{1.7cm}}
\toprule
N.º & Proposição & Experimento e o que foi medido & Fonte & Status \\
\midrule\endfirsthead
\toprule
N.º & Proposição & Experimento e o que foi medido & Fonte & Status \\
\midrule\endhead
1 & O neurônio é um circuito RC com limiar e reset~\eqref{eq:glutneuron} & No corpo de neurônios piramidais do córtex de rato (fatias) injetou-se uma corrente ruidosa, semelhante ao fluxo sináptico no cérebro vivo. A dependência da frequência média em relação à média e à dispersão da corrente é descrita pelo neurônio integrador com vazamento, se for acrescentada uma adaptação lenta da frequência. As faixas de $\tau$ e dos atrasos são dadas no capítulo sem referência experimental & \cite{Rauch2003} & medido \\
2 & O modelo~\eqref{eq:glutneuron} é uma simplificação: os ramos de entrada emitem eles mesmos disparos locais & Registro direto dos dendritos de neurônios piramidais do córtex visual do camundongo in vivo: o estímulo visual evoca disparos dendríticos locais que dependem dos receptores NMDA; suprimi-los alarga a sintonia do neurônio à orientação & \cite{Smith2013} & medido \\
3 & Janela de coincidência $\Delta_{\max}=\tau\ln\frac{w}{\theta-w}$~\eqref{eq:window}; com $\theta=1.5w$ ela vale $0.7\tau$ & Consequência do modelo~\eqref{eq:glutneuron}. Há medição direta de uma janela estreita de soma em neurônios com inibição rápida (capítulo~\ref{sec:gaba}, linha~3 de sua tabela) & dedução no capítulo & teoria \\
4 & Uma rede só de neurônios \nt{GLUT} calcula apenas funções monótonas das entradas (proposição~\ref{prop:monotone}) & Demonstração no capítulo: iteração a partir do silêncio com lado direito não decrescente. É uma afirmação sobre o modelo; não há experimento que a teste numa rede viva sem inibição & dedução no capítulo & teoria \\
5 & Os órgãos dos sentidos fornecem um par de fios: umas células respondem ao aumento do estímulo, outras à diminuição & Retina: já nas células bipolares, o sinal dos cones se divide em canais liga e desliga. Bloqueio farmacológico das bipolares liga no macaco: os canais seguem separados até o córtex; após o bloqueio piora a detecção do aumento de luz, mas não da diminuição & \cite{Schiller1992} & medido \\
6 & Com código de dois fios, uma rede de neurônios \nt{GLUT} calcula qualquer função lógica de modo assíncrono, sem ciclo comum, ao preço do dobro de neurônios (proposição~\ref{prop:dualrail}, \eqref{eq:dualrail}, fig.~\ref{fig:dualxor}) & O código de dois fios e a detecção de prontidão pelo próprio sinal são técnica padrão de circuitos assíncronos insensíveis a atrasos. O circuito de OU exclusivo com seis neurônios é montado no capítulo. Não há experimento mostrando que o cérebro calcula funções lógicas justamente com código de dois fios & \cite{Sparso2001} & teoria \\
7 & A maior diferença de tempo de chegada do som aos ouvidos no ser humano é $\approx0.6\ms$, e o cérebro distingue cerca de dez microssegundos & Ouvintes num experimento com duas alternativas de resposta: o limiar de discriminação da diferença de tempo (75\,\% de acertos) é de $9$\,\textmu s para ruído na banda de $150$--$1700$\,Hz, $11$\,\textmu s para um tom de $1$\,kHz, $28$\,\textmu s para um clique. O limite de $0.6\ms$ é cálculo do capítulo, $0.22\,\text{m}/343\,\text{m/s}$ & \cite{KlumppEady1956} & medido \\
8 & Nas aves, a diferença de tempos vira lugar por meio de linhas de atraso e detectores de coincidência~\eqref{eq:itd} (fig.~\ref{fig:itd}) & Suindara: os axônios dos dois ouvidos entram no núcleo dos detectores de coincidência por lados opostos; o tempo de chegada dos disparos ao longo do axônio muda de forma ordenada com a profundidade, e a variação dos atrasos através do núcleo ($160$\,\textmu s em $700$\,\textmu m) coincide com a faixa de atrasos interaurais & \cite{Carr1990} & medido \\
9 & Nos mamíferos, a fileira de atrasos não foi encontrada; a mesma tarefa é resolvida por detectores de coincidência com inibição de tempo preciso vinda de neurônios \nt{GLYC} & Registro in vivo na oliva superior medial do gerbil: as respostas não formam um mapa de direções; a aplicação de glicina e de seu bloqueador por micropipeta mostra que uma inibição \emph{glicinérgica} com tempo calculado com precisão define a faixa de trabalho dos atrasos. A inibição aqui vem de \nt{GLYC}, não de \nt{GABA} & \cite{Brand2002,Grothe2010} & medido \\
10 & Um grupo com forte realimentação é um flip-flop; a atividade autossustentada dura segundos após o estímulo (fig.~\ref{fig:bistable}) & Córtex pré-frontal do macaco numa tarefa de movimento ocular adiado para um ponto memorizado (atraso de $1$--$6$\,s): $87$ de $170$ neurônios ligados à tarefa mudam de frequência durante o atraso, e em $79\,\%$ deles isso depende da direção do ponto memorizado. Que essa atividade é mantida por uma realimentação com dois estados estáveis é teoria & \cite{Funahashi1989,Wang2001} & medido \\
11 & O bit é gravado se o aporte durar mais de $\approx0.7\tau$ (fig.~\ref{fig:recurrentgroup}b) & Cálculo da equação $\tau\,\dd r/\dd t=-r+f(wr+I)$ pelo método de Euler com $w=1$, $I=0.6$; não há script em \texttt{models/}. Não há experimento & cálculo no capítulo & modelo \\
12 & A memória pode ser guardada na facilitação das sinapses, quase sem disparos & Teoria: o cálcio residual nos terminais mantém o registro por cerca de um segundo. Base: no córtex pré-frontal medial do furão (registro de vários neurônios ao mesmo tempo) há uma sub-rede de neurônios piramidais ligados por sinapses facilitadoras com efeito prolongado. Revisão das observações de intervalos ``silenciosos'' durante a retenção na memória. Qual maneira o córtex usa e quando é questão em aberto & \cite{Mongillo2008,WangMarkram2006,Stokes2015} & hipótese \\
13 & A memória é apagada pela fadiga: o estoque de neurotransmissor se esgota & Registros pareados de neurônios piramidais do córtex: a resposta da sinapse cai com disparos frequentes, e a velocidade da queda é dada pela probabilidade de liberação. Que é exatamente isso que apaga o grupo autossustentado não foi mostrado diretamente & \cite{Tsodyks1997} & medido \\
14 & Uma cadeia de grupos é um temporizador disparado por evento; nas aves canoras, os neurônios disparam estritamente em sequência & Registro de neurônios identificados do centro vocal superior do mandarim durante o sono e o canto: cada neurônio que envia saída ao núcleo motor emite uma única rajada curta num momento preciso do canto, e os neurônios disparam em sequência & \cite{Hahnloser2002} & medido \\
\bottomrule
\end{longtable}}
